%%=============================================================================
%% Inleiding
%%=============================================================================

\chapter{\IfLanguageName{dutch}{Inleiding}{Introduction}}%
\label{ch:inleiding}

In de huidige gedigitaliseerde economie is de constante beschikbaarheid van vitale webapplicaties een fundamentele vereiste voor de continuïteit van organisaties. Bedrijven vertrouwen tegenwoordig volledig op hun IT-infrastructuur voor hun dagelijkse werking, communicatie en dienstverlening. Zelfs een verstoring van slechts enkele minuten heeft merkbare gevolgen binnen de bedrijfsvoering. Bovendien heeft downtime een domino-effect op de gehele organisatie, ongeacht of er gegevensverlies, datamutaties of beveiligingsincidenten plaatsvinden \autocite{ITIC2024}.
\medskip

Ongeplande downtime als gevolg van serveruitval leidt niet alleen tot grootschalige verliezen, maar tast ook de reputatie van een organisatie aan. Volgens recente rapportages van \textcite{ITIC2024} kunnen de kosten van downtime voor grote ondernemingen oplopen tot ver boven de \$300.000 per uur, waarbij factoren zoals verloren productiviteit en juridische claims de rekening verder doen stijgen. Naast deze directe financiële schade wijst onderzoek van \textcite{Kollective2025} op zogenoemde ‘verborgen kosten’, zoals een dalende werknemersmoraal en verlies van klantvertrouwen. Deze effecten zijn vaak moeilijker te kwantificeren, maar kunnen op de lange termijn minstens zo schadelijk zijn.
\medskip

Om deze risico's te beperken, is de implementatie van een High Availability (HA) cluster essentieel. HA verwijst naar het ontwerp van systemen die consistent beschikbaar blijven door de toepassing van drie kernprincipes: redundantie, veerkracht en failover. In dit onderzoek wordt specifiek gekozen voor een Active-Passive architectuur. De voordelen van dit model zijn de eenvoud in beheer en de kosteneffectiviteit voor applicaties met een gemiddeld verkeersvolume \autocite{Luca2024}.
\medskip

Voor de realisatie van de HA-cluster wordt gebruikgemaakt van open-source componenten op Red Hat Enterprise Linux (RHEL)-compatibele nodes. Corosync fungeert hierbij als de communicatielaag, terwijl Pacemaker de resource manager is. Onderzoek toont aan dat een omgeving beheerd door Pacemaker de downtime tot wel 45\% kan reduceren in vergelijking met niet-geoptimaliseerde systemen. Door de integratie van Pacemaker wordt de systeemveerkracht aanzienlijk verbeterd en blijft de verwerkingscapaciteit stabiel, zelfs tijdens zware werklasten \autocite{Mani2022}.

\section{\IfLanguageName{dutch}{Probleemstelling}{Problem Statement}}%
\label{sec:probleemstelling}

Hoewel de genoemde tools een robuuste basis leggen, is de standaardconfiguratie vaak niet toereikend om te voldoen aan strikte eisen voor de Recovery Time Objective (RTO). De RTO definieert de maximaal acceptabele tijd waarin een applicatie na een storing weer operationeel moet zijn. Voor veel moderne diensten is een hersteltijd van enkele seconden cruciaal. De grootste uitdaging ligt in het nauwkeurig afstellen van parameters zoals timeoutwaarden, resource stickiness en fencingmechanismen zoals STONITH. Zonder een effectief fencingmechanisme ontstaat het risico op 'split-brain'-scenario's, waarbij de data-integriteit niet langer gegarandeerd kan worden.
\medskip

Het onderzoek richt zich op IT-architecten, systeembeheerders en DevOps-engineers die verantwoordelijk zijn voor de implementatie en het beheer van Linux-gebaseerde infrastructuur en op zoek zijn naar betaalbare, open-source alternatieven voor commerciële HA-oplossingen.

\section{\IfLanguageName{dutch}{Onderzoeksvraag}{Research question}}%
\label{sec:onderzoeksvraag}

De centrale onderzoeksvraag van deze bachelorproef luidt:
\begin{quote}
    \textit{“Hoe kan een betaalbare Active-Passive High Availability webomgeving, bestaande uit Pacemaker en Corosync op RHEL-compatibele nodes, geoptimaliseerd worden om de Recovery Time Objective (RTO) van de webapplicatie bij serveruitval te minimaliseren?”}
\end{quote}
\medskip 

Om een onderbouwd antwoord op de centrale onderzoeksvraag te formuleren, wordt het onderzoek opgesplitst in een reeks gerichte deelvragen. In de eerste plaats wordt het probleemdomein in kaart gebracht om de theoretische en technische aspecten van de cluster-architectuur te doorgronden:
\begin{itemize}
    \item Wat zijn de specifieke functies van Pacemaker, Corosync en STONITH (fencing) binnen de gekozen HA-architectuur, en hoe werken deze drie componenten samen om een consensus over de status van de webdienst te bereiken?
    \item Welke Pacemaker-configuratieparameters en constraints (zoals op-start-interval, resource-stickiness en timeout waarden) beïnvloeden het meest direct de snelheid van de failover en daarmee de RTO?
    \item Welke vereisten en risico’s zijn cruciaal bij het gebruik van een virtueel IP-adres voor de webdienst in een Active-Passive cluster, en hoe beheert Pacemaker de migratie van dit virtuele IP om zero downtime bij failover te benaderen?
\end{itemize}
\medskip 

Voortbouwend op deze technische basis verschuift de focus naar het oplossingsdomein, waarbij de nadruk ligt op de effectieve implementatie en de kwantitatieve optimalisatie van de hersteltijd:
\begin{itemize}
    \item Hoe kan een robuust en betaalbaar Power Fencing-mechanisme (STONITH) - cruciaal voor het garanderen van data-integriteit - effectief geïmplementeerd en getest worden met behulp van een geschikte fence agent om ’split-brain’-scenario’s betrouwbaar te voorkomen?
    \item Wat is de gemeten gemiddelde Recovery Time Objective (RTO) van de webdienst, gesimuleerd voor ten minste drie verschillende uitvalscenario’s (bijvoorbeeld een harde nodecrash, een service crash, en netwerkisolatie)?
    \item Hoe kan de initiële Pacemaker-configuratie, op basis van de theoretisch geïdentificeerde parameters, stapsgewijs geoptimaliseerd worden om de RTO tot het laagst mogelijke niveau te reduceren, en wat is de impact van deze optimalisatie op de normale performantie van de webdienst?
\end{itemize}

\section{\IfLanguageName{dutch}{Onderzoeksdoelstelling}{Research objective}}%
\label{sec:onderzoeksdoelstelling}

Dit onderzoek richt zich op het identificeren en optimaliseren van de configuratie-parameters binnen een Active-Passive HA-webomgeving om de RTO bij serveruitval te minimaliseren. Het uiteindelijke resultaat is een functionele Proof-of-Concept (PoC) en een best-practice handleiding die als blauwdruk dient voor de doelgroep.

\section{\IfLanguageName{dutch}{Opzet van deze bachelorproef}{Structure of this bachelor thesis}}%
\label{sec:opzet-bachelorproef}

De rest van deze bachelorproef is als volgt opgebouwd:
\medskip 

In Hoofdstuk~\ref{ch:stand-van-zaken} worden de theoretische deelvragen uit het probleemdomein behandeld en wordt de technologische basis van de gebruikte tools dieper geanalyseerd. 

In Hoofdstuk~\ref{ch:methodologie} wordt de opzet van de testomgeving en de wijze waarop de RTO-metingen worden uitgevoerd toegelicht. 

In Hoofdstuk~\ref{ch:conclusie} wordt een finaal antwoord geformuleerd op de onderzoeksvraag en worden aanbevelingen gedaan voor de praktijk.